\Article{50}
Let $X$ denote a function of the indeterminate $x$ of the following form:
\[
Ax^n+Bx^{n-1}+Cx^{n-2}+\text{etc.}+Mx+N;
\]

\SourceLeaf{131}%
where $n$ is a positive real integer, $A,B,C$, etc., real or imaginary integers, and $m$ a complex integer. Here too we shall call any integer a \emph{root} of the congruence $X\equiv0\pmod m$ if, when substituted for $x$, it gives $X$ a value divisible by the modulus $m$. We shall not regard solutions given by roots congruent modulo the modulus as different.

Whenever the modulus is a prime number, here too such a congruence of degree $n$ cannot have more than $n$ different solutions. Let $\alpha$ denote any fixed integer (complex). $X$ can be put by division by $x-\alpha$, with the indeterminate left free, in the form $X=(x-\alpha)X'+h$, where $h$ is a fixed integer and $X'$ a function of degree $n-1$ with integer coefficients. Whenever $\alpha$ is a root of the congruence $X\equiv0\pmod m$, $h$ is evidently divisible by $m$; equivalently, with the indeterminate left free, we have $X\equiv(x-\alpha)X'\pmod m$.

Similarly, if $\beta$ denotes a fixed integer and $X'$ is put in the form $(x-\beta)X''+h'$, $X''$ will be a function of degree $n-2$ with integer coefficients. If $\beta$ is supposed to be a root of the congruence $X\equiv0$, it must also satisfy $(\beta-\alpha)X'\equiv0$, and also $X'\equiv0$ if the roots $\alpha,\beta$ are incongruent. We conclude that $h'$ must also be divisible by $m$; equivalently, with the indeterminate left free, $X\equiv(x-\alpha)(x-\beta)X''\pmod m$.

\Needspace{7\baselineskip}
In the same way, on adding a third root $\gamma$, incongruent to the preceding ones, we shall have, with the indeterminate left free, $X\equiv(x-\alpha)(x-\beta)(x-\gamma)X'''$, where $X'''$ is a function of degree $n-3$ with integer coefficients. We may continue in the same way. It is also clear that the coefficient of the highest term in each function is $=A$, which we may suppose not divisible by $m$, since otherwise the congruence $X\equiv0$ would essentially have to be assigned a lower degree. Thus, whenever there are $n$ incongruent roots, say $\alpha,\beta,\gamma\dotsfour\nu$, we have, with the indeterminate left free,
\[
X\equiv A(x-\alpha)(x-\beta)(x-\gamma)\dotsfour(x-\nu)\pmod m
\]
Consequently, substituting a new value incongruent to each of $\alpha,\beta,\gamma\dotsfour\nu$ would certainly give $X$ a value not divisible by $m$, from which the theorem follows at once.

This proof essentially agrees with the one we gave in \emph{Disq. Ar.} Article~43, every step of which applies to complex numbers just as it does to real numbers.

\SourceLeaf{132}%
\Article{51}
Most of what was given in the third section of the \emph{Disquisitiones Arithmeticae} concerning residues of powers also holds in the arithmetic of complex numbers, with slight changes; indeed, the proofs of the theorems could generally be retained. So that nothing is missing, however, we shall present the principal theorems supported by concise proofs. Throughout, the modulus is to be understood to be a prime number.

\textsc{Theorem.} \emph{If $k$ denotes an integer not divisible by the modulus $m$, of norm $=p$, then $k^{p-1}\equiv1\pmod m$.}

\emph{Proof.} Let $a,b,c$, etc., form a complete system of incongruent residues modulo $m$, but with the residue divisible by $m$ omitted, so that the number of these numbers, whose collection we denote by $C$, is $=p-1$. Further, let $C'$ be the collection of products $ka,kb,kc$, etc. By hypothesis, none of these products is divisible by $m$, so each has a congruent residue in the collection $C$: we may thus have $ak\equiv a',bk\equiv b',ck\equiv c'$, etc., $(\operatorname{mod.}m)$, where the numbers $a',b',c'$, etc., themselves occur in the collection $C$. Denote the collection of $a',b',c'$, etc., by $C''$. Let $P,P',P''$ be the products of the individual numbers in the collections $C,C',C''$, respectively; that is,
\[
\begin{aligned}
P&=abc\dotsfour\\
P'&=k^{p-1}abc\dotsfour=k^{p-1}P\\
P''&=a'b'c'\dotsfour
\end{aligned}
\]
\Needspace{9\baselineskip}
Since the numbers in $C''$ are congruent, in succession, to the numbers in $C'$, we have $P''\equiv P'$, or $P''\equiv k^{p-1}P$. But any two numbers in $C''$ are readily seen to be incongruent, and therefore all are distinct. The numbers in $C''$ must consequently coincide exactly with those in $C$, with only their order changed, giving $P''=P$. Thus $(k^{p-1}-1)P$ is divisible by $m$. Since $m$ is a prime number dividing none of the individual factors of $P$, $k^{p-1}-1$ must necessarily be divisible by $m$. Q.~E.~D.

\Article{52}
\textsc{Theorem.} \emph{If $k$, as in the preceding article, denotes an integer not divisible by the modulus $m$, and $t$ the least exponent (apart from $0$) for which $k^t\equiv1\pmod m$, then $t$ divides every other exponent $u$ for which $k^u\equiv1\pmod m$.}

\SourceLeaf{133}%
\emph{Proof.} If $t$ did not divide $u$, let $gt$ be the multiple of $u$ next greater than $u$, so that $gt-u$ is a positive integer less than $t$. From $k^t\equiv1,k^u\equiv1$ it follows that $0\equiv k^{gt}-k^u\equiv k^u(k^{gt-u}-1)$, and hence $k^{gt-u}\equiv1$: there is a power of $k$ with an exponent less than $t$ that is congruent to unity, contrary to the hypothesis.

It follows as a corollary that $t$ certainly divides $p-1$.

Here too we shall call numbers $k$ for which $t=p-1$ \emph{primitive roots} for. the modulus $m$. We shall now show that such numbers really exist.

\Article{53}
Factor the number $p-1$ into its prime factors, so that
\[
p-1=a^\alpha b^\beta c^\gamma\dotsfour
\]
where $a,b,c$, etc., denote distinct positive real prime numbers. Let $A,B,C$, etc., be integers (complex) not divisible by $m$, and respectively \emph{not} satisfying the congruences
\[
x^{\frac{p-1}a}\equiv1,\quad x^{\frac{p-1}b}\equiv1,\quad x^{\frac{p-1}c}\equiv1\quad\text{etc.}
\]
modulo $m$. The existence of such numbers is clear from the theorem of Article~50. Finally, let $h$ be congruent modulo $m$ to the product
\[
A^{\frac{p-1}{a^\alpha}}B^{\frac{p-1}{b^\beta}}C^{\frac{p-1}{c^\gamma}}\dotsfour
\]
Then I say that $h$ will be a primitive root.

\emph{Proof.} Denoting by $t$ the exponent of the lowest power $h^t$ congruent to unity, if $h$ were not a primitive root, $t$ would be a proper divisor of $p-1$; that is, $\frac{p-1}t$ would be an integer greater than unity. This integer evidently has its real prime factors among $a,b,c$, etc. Let us therefore suppose, as we may, that $\frac{p-1}t$ is divisible by $a$, and put $p-1=atu$. Since $h^t\equiv1$, we then also have $h^{tu}\equiv1$, or
\[
A^{\frac{p-1}{a^\alpha}\pdot\frac{p-1}a}B^{\frac{p-1}{b^\beta}\pdot\frac{p-1}a}C^{\frac{p-1}{c^\gamma}\pdot\frac{p-1}a}\dotsfour\equiv1
\]
But $\frac{p-1}{ab^\beta}$ is evidently an integer, so
\[
B^{\frac{p-1}{b^\beta}\pdot\frac{p-1}a}=(B^{p-1})^{\frac{p-1}{ab^\beta}}\equiv1
\]
and similarly
\SourceLeaf{134}%
$C^{\frac{p-1}{c^\gamma}\pdot\frac{p-1}a}\equiv1$, and so on; consequently we must have \mbox{$A^{\frac{p-1}{a^\alpha}\pdot\frac{p-1}a}\equiv1$}
Now determine a positive integer $\lambda$ such that
\[
\lambda b^\beta c^\gamma\dotsfour\equiv1\pmod a
\]
This can be done, since the prime number $a$ does not divide $b^\beta c^\gamma\dotsfour$; put $\lambda b^\beta c^\gamma\dotsfour=1+a\mu$. Evidently $A^{\lambda\pdot\frac{p-1}{a^\alpha}\pdot\frac{p-1}a}\equiv1$; or, since $\lambda\pdot\frac{p-1}{a^\alpha}\pdot\frac{p-1}a=(1+a\mu)\frac{p-1}a=(p-1)\mu+\frac{p-1}a$, we have $A^{(p-1)\mu}\pdot A^{\frac{p-1}a}\equiv1$. Hence, since $A^{(p-1)\mu}\equiv1$ holds of itself, we also have $A^{\frac{p-1}a}\equiv1$, contrary to the hypothesis. The supposition that $t$ is a proper divisor of $p-1$ therefore cannot stand, and $h$ must necessarily be a primitive root.

